%% LyX 2.3.3 created this file.  For more info, see http://www.lyx.org/.
%% Do not edit unless you really know what you are doing.
\documentclass[english]{article}
\usepackage[T1]{fontenc}
\usepackage[latin9]{inputenc}
\usepackage{geometry}
\geometry{verbose,tmargin=1in,bmargin=1in,lmargin=1in,rmargin=1in}
\usepackage{mathtools}
\usepackage{bm}
\usepackage{amsmath}
\usepackage{amssymb}

\makeatletter
%%%%%%%%%%%%%%%%%%%%%%%%%%%%%% User specified LaTeX commands.


\usepackage{amsthm} % for theorems

\usepackage{hyperref} % hyperrefs


\usepackage{amsmath,amssymb} % for \vertiii
\newcommand{\vertiii}[1]{{\left\vert\kern-0.25ex\left\vert\kern-0.25ex\left\vert #1 
    \right\vert\kern-0.25ex\right\vert\kern-0.25ex\right\vert}}

\usepackage{bbm}


\let\hat\widehat
\let\tilde\widetilde
\let\check\widecheck
\let\epsilon\varepsilon



\usepackage{color}
\definecolor{cm}{RGB}{0,0,200}
\newcommand{\cm}[1]{\textcolor{cm}{[CM: #1]}}

\makeatother

\usepackage{babel}
\begin{document}
\theoremstyle{plain} \newtheorem{lemma}{\textbf{Lemma}} \newtheorem{prop}{\textbf{Proposition}}\newtheorem{theorem}{\textbf{Theorem}}\setcounter{theorem}{0}
\newtheorem{corollary}{\textbf{Corollary}} \newtheorem{assumption}{\textbf{Assumption}}
\newtheorem{example}{\textbf{Example}} \newtheorem{definition}{\textbf{Definition}}
\newtheorem{fact}{\textbf{Fact}} \theoremstyle{definition}

\theoremstyle{remark}\newtheorem{remark}{\textbf{Remark}}\newtheorem{condition}{Condition}\newtheorem{claim}{\textbf{Claim}}
\title{Classical perturbation theory}

\maketitle
\tableofcontents{}

\vspace{2.5em}

Let $\bm{A}\in\mathbb{R}^{m\times n}$ be the matrix under study and
$\hat{\bm{A}}=\bm{A}+\bm{E}$ be the perturbed matrix, where $\bm{E}\in\mathbb{R}^{m\times n}$
represents the perturbation. The central question of the current monograph
is: How does the eigenvalues (resp.~singular values) and eigenvectors
(resp.~singular vectors) of $\bm{A}$ change in response to the perturbation~$\bm{E}$?

\section{Mathematical backgrounds}

Before presenting the classical matrix perturbation theory, we first
review some basic facts in matrix analysis and probability theory
that are used extensively in the coming sections. The readers familiar
with these materials can proceed directly to Section~\ref{sec:classical-perturbation}.

\subsection{Matrix analysis \label{subsec:Matrix-analysis}}

\subsubsection{Eigenvalues and eigenvectors}

We start with the definition of eigenvalues and eigenvectors of a
square matrix.

\begin{definition}[Eigenpair]\label{def:eigenpair} Fix a square
matrix $\bm{A}\in\mathbb{C}^{n\times n}$. We say $(\lambda,\bm{v})\in\mathbb{C}\times\mathbb{C}^{n}$
is an \emph{eigenpair} of $\bm{A}$ if $\bm{v}$ is nonzero and 
\begin{equation}
\bm{A}\bm{v}=\lambda\bm{v}.\label{eq:def-eigenpair}
\end{equation}
In particular, the scalar $\lambda$ is an eigenvalue of $\bm{A}$
and $\bm{v}$ is the associated eigenvector. \end{definition}

In words, an eigenvector $\bm{v}$ of $\bm{A}$ does not change its
direction when multiplied by $\bm{A}$ and the eigenvalue $\lambda$
measures the amount of expansion or shrinkage in this transformation.
Several remarks regarding Definition~\ref{def:eigenpair} are in
order.
\begin{itemize}
\item The eigenpair (i.e.~the eigenvalue and the eigenvector) is only defined
for square matrices. This is implicitly assumed in the definition
of the eigenpair~(see~(\ref{def:eigenpair})).
\item Every $n\times n$ square matrix has exactly $n$ (complex) eigenvalues,
possibly with repetitions. To see this, one can transform (\ref{def:eigenpair})
into the identity
\[
\left(\lambda\bm{I}_{n}-\bm{A}\right)\bm{v}=\bm{0}.
\]
Consequently, the scalar $\lambda$ is an eigenvalue of $\bm{A}$
if and only if 
\[
\mathsf{det}\left(\lambda\bm{I}_{n}-\bm{A}\right)=0.
\]
Notice that the above relation specifies a polynomial equation of
degree $n$, which by the fundamental theorem of algebra, has exactly
$n$ roots in the complex field (and hence $n$ complex eigenvalues
of $\bm{A}$).
\item A matrix consisting of real entries can have complex eigenvalues and
complex eigenvectors. Take for instance the matrix
\[
\bm{A}=\left[\begin{array}{cc}
3 & -2\\
4 & -1
\end{array}\right],
\]
where $1+2i$ and $1-2i$ are its two complex eigenvalues. Here $i$
denotes the imaginary unit.
\end{itemize}
With the last remark being said, there exists one important set of
matrices --- \emph{real symmetric matrices} --- of which the eigenvalues
must be real. This claim is formalized in the following theorem.

\begin{theorem}[Spectral theorem]\label{thm:spectral} Let $\bm{A}\in\mathbb{R}^{n\times n}$
be a real symmetric matrix. Then there exist $n$ real eigenvalues
$\lambda_{1},\cdots\lambda_{n}\in\mathbb{R}$ and an orthonormal basis
$\{\bm{u}_{1},\bm{u}_{2},\cdots,\bm{u}_{n}\}\subset\mathbb{R}^{n}$
such that $\bm{A}\bm{u}_{i}=\lambda_{i}\bm{u}_{i}$. Equivalently,
a real symmetric matrix $\bm{A}$ enjoys the following spectral decomposition
\begin{equation}
\bm{A}=\sum_{i=1}^{n}\lambda_{i}\bm{u}_{i}\bm{u}_{i}^{\top},\label{eq:spectral-decomposition-summation-form}
\end{equation}
where $\{\lambda_{i}\}_{1\leq i\leq n}$ are the real eigenvalues
of $\bm{A}$ and $\{\bm{u}_{i}\}_{1\leq i\leq n}$ is an orthonormal
basis of~$\mathbb{R}^{n}$.

\end{theorem}

Since the eigenvalues of real symmetric matrices are guaranteed to
be real numbers, we can therefore order them. From now on, denote
by $\lambda_{i}(\bm{A})$ the $i$th largest eigenvalue of a real
symmetric matrix $\bm{A}\in\mathbb{R}^{n\times n}$. In addition,
we define $\lambda_{\max}(\bm{A})\triangleq\lambda_{1}(\bm{A})$ and
$\lambda_{\min}(\bm{A})\triangleq\lambda_{n}(\bm{A})$.

To make the expression~(\ref{eq:spectral-decomposition-diagonalization-form})
more compact, we denote 
\[
\bm{U}=\left[\begin{array}{ccc}
\rvert &  & \rvert\\
\bm{u}_{1} & \cdots & \bm{u}_{n}\\
\rvert &  & \rvert
\end{array}\right]\in\mathbb{R}^{n\times n}\qquad\text{and}\qquad\bm{\Lambda}=\left[\begin{array}{cccc}
\lambda_{1} & 0 & \cdots & 0\\
0 & \lambda_{2} & \cdots & 0\\
\vdots & \cdots & \ddots & \vdots\\
0 & 0 & \cdots & \lambda_{n}
\end{array}\right]\in\mathbb{R}^{n\times n}.
\]
As a result, the decomposition (\ref{eq:spectral-decomposition-diagonalization-form})
can be rewritten as 
\begin{equation}
\bm{A}=\bm{U}\bm{\Lambda}\bm{U}^{\top},\label{eq:spectral-decomposition-multiplication-form}
\end{equation}
in which a real symmetric matrix is decomposed into the multiplications
of orthonormal matrices (i.e.~$\bm{U}$) and a real diagonal matrix
(i.e.~$\bm{\Lambda}$). Another useful way to interpret spectral
decomposition is through the lens of diagonalization. Simple algebraic
manipulation reveals an equivalent form of (\ref{eq:spectral-decomposition-multiplication-form})
\begin{equation}
\bm{U}^{\top}\bm{A}\bm{U}=\bm{\Lambda}.\label{eq:spectral-decomposition-diagonalization-form}
\end{equation}
This essentially tells us that any real symmetric matrix $\bm{A}$
can be \emph{diagonalized }by an orthonormal matrix~(i.e.~$\bm{U}$).

\subsubsection{Singular values and singular vectors}

As we have pointed out in the remarks following Definition~\ref{def:eigenpair},
a rectangular matrix does \emph{not} have well-defined eigenvalues
and eigenvectors. It is then natural to ask whether there is a similar
decomposition of a general matrix to that of a real symmetric matrix
(cf.~(\ref{eq:spectral-decomposition-multiplication-form})). Fortunately,
a general matrix admits another useful representation, namely the
singular value decomposition (SVD).

\begin{definition}[Singular value decomposition]\label{def:SVD}Let
$\bm{A}\in\mathbb{R}^{m\times n}$ be a general matrix. Without loss
of generality, assume that $m\leq n$. There exists a pair of orthogonal
matrices $\bm{U}\in\mathbb{R}^{m\times m}$ and $\bm{V}\in\mathbb{R}^{n\times n}$
such that 
\begin{equation}
\bm{A}=\bm{U}\bm{\Sigma}\bm{V}^{\top},\label{eq:SVD}
\end{equation}
where $\bm{\Sigma}$ is a rectangular diagonal matrix with 
\[
\bm{\Sigma}=\left[\begin{array}{cccccc}
\sigma_{1} & 0 & \cdots & \cdots & 0 & 0\\
0 & \sigma_{2} & \cdots & \cdots & 0 & 0\\
\vdots & \vdots & \ddots & \vdots & \vdots & \vdots\\
0 & 0 & \cdots & \sigma_{m} & 0 & 0
\end{array}\right]\in\mathbb{R}^{m\times n}
\]
and $\sigma_{1}\geq\sigma_{2}\geq\cdots\geq\sigma_{m}\geq0$. The
decomposition (\ref{eq:SVD}) is called the singular value decomposition,
where the nonnegative real numbers $\{\sigma_{i}\}_{1\leq i\leq n}$
are the singular values of $\bm{A}$ and the columns in $\bm{U}$
(resp.~$\bm{V}$) are the left (resp.~right) singular vectors of
$\bm{A}$. \end{definition}

Similar to the eigenvalues of real symmetric matrices, the singular
values of a general matrix are guaranteed to be real. As a result,
for any $1\leq i\leq\min\{m,n\}$, we define $\sigma_{i}(\bm{A})$
to be $i$th largest singular value of $\bm{A}$, and define $\sigma_{\max}(\bm{A})\triangleq\sigma_{1}(\bm{A})$
(resp.~$\sigma_{\min}(\bm{A})\triangleq\sigma_{\min\{m,n\}}(\bm{A})$)
to be the largest (resp.~smallest) singular values of $\bm{A}$.

It is straightforward to read the rank of a matrix from its singular
values in that the rank is equal to the number of nonzero singular
values. If the rank $r$ of a matrix $\bm{A}\in\mathbb{R}^{m\times n}$
is smaller than $\min\{m,n\}$, it is often more economical to represent
the matrix using the \emph{compact SVD }
\begin{equation}
\bm{A}=\bm{U}_{r}\bm{\Sigma}_{r}\bm{V}_{r}^{\top},\label{eq:compact-SVD}
\end{equation}
where $\bm{U}_{r}\in\mathbb{R}^{m\times r}$ (resp.~$\bm{V}_{r}\in\mathbb{R}^{n\times r}$)
is the first $r$ columns of $\bm{U}$ (resp.~$\bm{V}$) and $\bm{\Sigma}_{r}\in\mathbb{R}^{r\times r}$
is the upper left block of $\bm{\Sigma}$ of size $r$.

Having introduced two matrix decompositions, it is instrumental to
compare the singular value decomposition of a general matrix (cf.~Definition~\ref{def:SVD})
with the spectral decomposition of a real symmetric matrix (cf.~Theorem~(\ref{thm:spectral})).
First, the eigenvalues of a real symmetric matrix are not necessarily
nonnegative whereas singular values must be. Second, as claimed in
(\ref{eq:spectral-decomposition-diagonalization-form}), any real
symmetric matrix can be diagonalized by a single orthonormal matrix.
However, this may not hold true for a general matrix, nevertheless,
SVD points out a way to diagonalize a matrix using two orthonormal
matrices. To see this, one can equate (\ref{eq:SVD}) with

\[
\bm{U}^{\top}\bm{A}\bm{V}=\bm{\Sigma}.
\]
Last but not least, the singular value decomposition and the spectral
decomposition are inherently related. It is an easy exercise to see
that the columns of $\bm{U}$ (resp.~$\bm{V}$) are eigenvectors
of $\bm{A}\bm{A}^{\top}$ (resp.~$\bm{A}^{\top}\bm{A}$). In addition,
the singular values of $\bm{A}$ are nonnegative square roots of the
eigenvalues of either $\bm{A}\bm{A}^{\top}$ or $\bm{A}^{\top}\bm{A}$
(modulo zero singular values). This connection shall prove useful
in later discussions.

\subsubsection{Matrix norms}

To measure the size of the perturbation (e.g.~$\bm{E}$), we need
to introduce appropriate metrics~/~norms on the space of matrices.
Among all matrix norms, unitarily invariant norms are of most interest
to us. \begin{definition}[Unitarily invariant norms]A matrix norm
$\vertiii{\cdot}$ on $\mathbb{R}^{m\times n}$ is unitarily invariant
if
\[
\vertiii{\bm{A}}=\vertiii{\bm{U}^{\top}\bm{A}\bm{V}}
\]
for any matrix $\bm{A}\in\mathbb{R}^{m\times n}$ and any two orthogonal
matrices $\bm{U}\in\mathbb{R}^{m\times m}$ and $\bm{V}\in\mathbb{R}^{n\times n}$.
\end{definition}

Two prominent examples include the Frobenius norm and the operator
norm

\begin{equation}
\left\Vert \bm{A}\right\Vert _{\mathrm{F}}\triangleq\sqrt{\sum_{i,j}A_{ij}^{2}}\qquad\text{and}\qquad\left\Vert \bm{A}\right\Vert \triangleq\max_{\bm{x}\in\mathbb{R}^{n}:\bm{x}\neq\bm{0}}\frac{\left\Vert \bm{A}\bm{x}\right\Vert _{2}}{\left\Vert \bm{x}\right\Vert _{2}}.\label{eq:def-fro-op-norm}
\end{equation}
To see why both $\|\cdot\|_{\mathrm{F}}$ and $\|\cdot\|$ are unitarily
invariant norms, we have the following important characterizations
\[
\left\Vert \bm{A}\right\Vert _{\mathrm{F}}=\sqrt{\sum_{i=1}^{\min\{m,n\}}\sigma_{i}^{2}\left(\bm{A}\right)}\qquad\text{and}\qquad\left\Vert \bm{A}\right\Vert =\sigma_{1}\left(\bm{A}\right),
\]
where we recall that $\sigma_{1}(\bm{A})\geq\sigma_{2}(\bm{A})\geq\cdots\geq\sigma_{\min\{m,n\}}(\bm{A})\geq0$
are the singular values of $\bm{A}$ in descending order.

Last but not least, unitarily invariant norms enjoy the following
nice properties.

\begin{prop}Let $\vertiii{\cdot}$ be any unitarily invariant norm.
One has 
\begin{align*}
\vertiii{\bm{A}\bm{B}} & \leq\vertiii{\bm{A}}\left\Vert \bm{B}\right\Vert _{2}\qquad\text{and}\qquad\vertiii{\bm{A}\bm{B}}\leq\vertiii{\bm{B}}\left\Vert \bm{A}\right\Vert _{2};\\
\vertiii{\bm{A}\bm{B}} & \geq\vertiii{\bm{A}}\sigma_{\min}\left(\bm{B}\right)\qquad\text{and}\qquad\vertiii{\bm{A}\bm{B}}\leq\vertiii{\bm{B}}\sigma_{\min}\left(\bm{A}\right).
\end{align*}
\end{prop}

\subsection{Concentration inequalities}

\cm{Need to motivate why we need these. }

As we have mentioned in Section~XXX, we often observe

One quantitative version of the central limit theorem \cm{TODO}

\begin{theorem}[Bernstein's inequality]Let $\{x_{i}\}_{1\leq i\leq n}$
be a sequence of independent random variables. Suppose that for each
$1\leq i\leq n$, $\mathbb{E}[x_{i}]=0$ and $|x_{i}|\leq B$ almost
surely. Then we have for all $t\geq0$ 
\[
\mathbb{P}\left(\left|\sum_{i=1}^{n}x_{i}\right|\geq t\right)\leq2\exp\left(\frac{-t^{2}/2}{V+Bt/3}\right),
\]
where $V$ is the variance statistic 
\[
V\triangleq\left|\sum_{i=1}^{n}\mathbb{E}\left[x_{i}^{2}\right]\right|.
\]
\end{theorem}

\begin{theorem}[Matrix Bernstein's inequality]Let $\{\bm{X}_{i}\}_{1\leq i\leq n}$
be a sequence of independent real random matrices with dimension $d_{1}\times d_{2}$.
Suppose that 
\[
\mathbb{E}\left[\bm{X}_{i}\right]=\bm{0},\qquad\left\Vert \bm{X}_{l}\right\Vert \leq B\qquad\text{for each }i.
\]
Let $V$ be the variance statistic
\[
V\triangleq\max\left\{ \left\Vert \sum_{i=1}^{n}\mathbb{E}\left[\bm{X}_{i}\bm{X}_{i}^{\top}\right]\right\Vert ,\left\Vert \sum_{i=1}^{n}\mathbb{E}\left[\bm{X}_{i}^{\top}\bm{X}_{i}\right]\right\Vert \right\} .
\]
Then one has for all $t\geq0$, 
\[
\mathbb{P}\left(\left\Vert \sum_{i=1}^{n}\bm{X}_{i}\right\Vert \geq t\right)\leq\left(d_{1}+d_{2}\right)\exp\left(\frac{-t^{2}/2}{V+Bt/3}\right).
\]
\end{theorem},

\section{The classical $\ell_{2}$ perturbation theory \label{sec:classical-perturbation}}

Now we are ready to present the classical matrix perturbation theory
on the eigenvalues (resp.~singular values) and eigenvectors (resp.~singular
vectors).

\subsection{Perturbation of real symmetric matrices \label{subsec:Perturbation-of-real-symmetric}}

We start with a real symmetric matrix $\bm{A}\in\mathbb{R}^{n\times n}$,
where the focus is its eigenvalues and eigenvectors. In light of Theorem~\ref{thm:spectral},
the matrix $\bm{A}$ has the spectral decomposition 
\[
\bm{A}=\sum_{i=1}^{n}\lambda_{i}\bm{u}_{i}\bm{u}_{i}^{\top},
\]
where $\lambda_{1}\geq\lambda_{2}\geq\cdots\geq\lambda_{n}$ are the
eigenvalues and $\{\bm{u}_{i}\}$ --- the set of eigenvectors --- forms
an orthonormal basis of $\mathbb{R}^{n}$. Suppose that $\bm{A}$
is perturbed by a symmetric matrix $\bm{E}\in\mathbb{R}^{n\times n}$
and we only get to observe $\hat{\bm{A}}=\bm{A}+\bm{E}.$ Two important
questions follow: \cm{Need to motivate eigenspaces.}
\begin{enumerate}
\item How do the eigenvalues change when the matrix $\bm{A}$ is perturbed
by $\bm{E}$?
\item How do the eigenspaces change in response to the perturbation $\bm{E}$?
\end{enumerate}
Before addressing these questions, we introduce some useful notation.
Denote the spectral decomposition of $\hat{\bm{A}}$ by 
\[
\hat{\bm{A}}=\bm{A}+\bm{E}=\sum_{i=1}^{n}\hat{\lambda}_{i}\hat{\bm{u}}_{i}\hat{\bm{u}}_{i}^{\top},
\]
where $\hat{\lambda}_{1}\geq\hat{\lambda}_{2}\geq\cdots\geq\hat{\lambda}_{n}$
are the eigenvalues of $\hat{\bm{A}}$ and $\{\hat{\bm{u}}_{i}\}$
are the corresponding eigenvectors that are orthonormal.

\subsubsection{The perturbation of eigenvalues}

With all the notation at hand, we are in position to present the perturbation
bounds for eigenvalues of symmetric matrices.

\begin{lemma}[Weyls's inequality]\label{lemma:weyl}Instate the
notation in Section~\ref{subsec:Perturbation-of-real-symmetric}.
For every $1\leq i\leq n$, we have 
\begin{equation}
\left|\lambda_{i}\left(\bm{A}+\bm{E}\right)-\lambda_{i}\left(\bm{A}\right)\right|=\left|\hat{\lambda}_{i}-\lambda_{i}\right|\leq\left\Vert \bm{E}\right\Vert .\label{eq:weyl}
\end{equation}
\end{lemma}

An immediate implication of Lemma~\ref{lemma:weyl} is that the eigenvalues
of a real symmetric matrix are stable against small perturbations.

\subsubsection{Distance between subspaces}

Moving on to the question of the perturbation of eigenvectors~/~eigenspaces,
one clearly needs a distance metric on the space of linear subspaces.
Let $\mathcal{X}$ and $\mathcal{Y}$ be two linear subspaces of $\mathbb{R}^{n}$
with dimension $r$ and let $\bm{U}_{0}\in\mathbb{R}^{n\times r}$
(resp.~$\hat{\bm{U}}_{0}\in\mathbb{R}^{n\times r}$) be an orthonormal
basis of $\mathcal{X}$ (resp.~$\mathcal{Y}$). This subsection is
centered around quantifying the distance between $\mathcal{X}$ and
$\mathcal{Y}$. One naive idea is to use the ``metric'' $\vertiii{\bm{U}_{0}-\hat{\bm{U}}_{0}}$,
where $\vertiii{\cdot}$ is any norm on $\mathbb{R}^{n\times r}$.
However, an obvious drawback of this approach is that it does not
take into account the rotation ambiguity that is inherent in $\bm{U}_{0}$
and $\hat{\bm{U}}_{0}$ --- for any $\bm{R}\in\mathcal{O}^{r\times r}$,
the matrix $\bm{U}\bm{R}$ is also an orthonormal basis of $\mathcal{X}$.
Therefore, any reasonable distance metric should handle the rotation
ambiguity properly. Below are two popular metrics to measure the distance
between a pair of subspaces.
\begin{enumerate}
\item \emph{The algebraic construction of distance metrics --- projection
matrices. }It is an established fact in linear algebra that the projection
matrix $\bm{U}_{0}\bm{U}_{0}^{\top}$ (resp.~$\hat{\bm{U}}_{0}\hat{\bm{U}}_{0}^{\top}$)
onto the subspace $\mathcal{X}$ (resp.~$\mathcal{Y}$) is unique.
Utilizing this fact, one can define the distance between $\mathcal{X}$
and $\mathcal{Y}$ to be 
\begin{equation}
\rho_{\vertiii{\cdot}}\left(\mathcal{X},\mathcal{Y}\right)\triangleq\vertiii{\bm{U}_{0}\bm{U}_{0}^{\top}-\hat{\bm{U}}_{0}\hat{\bm{U}}_{0}^{\top}},\label{eq:algebraic-dist}
\end{equation}
where $\vertiii{\cdot}$ is some norm on $\mathbb{R}^{n\times n}$.
It is an exercise for the readers to prove that this is indeed a valid
distance metric on the space of linear subspaces.
\begin{enumerate}
\item \emph{The geometric construction of distance metrics --- canonical
angles}. Let $\bm{X}\bm{\Sigma}\bm{Y}^{\top}$ be the SVD of $\bm{U}_{0}^{\top}\hat{\bm{U}}_{0}$,
where $\bm{X},\bm{Y}\in\mathcal{O}^{r\times r}$ are orthonormal matrices
and $\bm{\Sigma}=\mathsf{diag}(\sigma_{1},\sigma_{2},\cdots,\sigma_{r})$
consisting of the singular values of $\bm{U}_{0}^{\top}\hat{\bm{U}}_{0}$
in a descending order. It is easy to check that all the singular values
$\{\sigma_{i}\}$ are contained in the interval $[0,1]$. Therefore,
one can define the \emph{canonical angles} between two subspaces spanned
by $\bm{U}_{0}$ (i.e.~$\mathcal{X}$) and $\hat{\bm{U}}_{0}$ (i.e.~$\mathcal{Y}$)
to be 
\[
\theta_{i}\triangleq\cos^{-1}\left(\sigma_{i}\right)\in\left[0,\pi/2\right],\qquad\text{for all}\quad1\leq i\leq r.
\]
Note that $0\leq\theta_{1}\leq\theta_{2}\leq\cdots\leq\theta_{r}\leq\pi/2$.
To see why this definition makes sense, let us consider the simplest
case where $r=1$ and hence $\bm{U}$ and $\hat{\bm{U}}$ are two
unit vectors. In this case, the canonical angle $\theta_{1}$ is exactly
the angle between the two unit vectors $\bm{U}$ and $\hat{\bm{U}}$:
when $\theta_{1}=0$, the two subspaces coincide with each other and
when $\theta_{1}=\pi/2$, they are far way from each other. Based
on this observation, the set of canonical angles $\{\theta_{i}\}$
forms a good basis to measure the distance between two subspaces.
To formalize this idea, we arrange the canonical angles into a diagonal
matrix 
\[
\bm{\Theta}\triangleq\mathsf{diag}\left(\theta_{1},\theta_{2},\cdots,\theta_{r}\right)\in\mathbb{R}^{r\times r},
\]
and correspondingly define 
\begin{align*}
\sin\bm{\Theta} & \triangleq\mathsf{diag}\left(\sin\theta_{1},\sin\theta_{2},\cdots,\sin\theta_{r}\right)\in\mathbb{R}^{r\times r}.
\end{align*}
With this notation in place, one can measure the distance between
$\mathcal{X}$ and $\mathcal{Y}$ through \cm{May not be a good notation
to use due to the ignorance of the norm. }
\begin{equation}
\rho_{\sin}\left(\mathcal{X},\mathcal{Y}\right)=\vertiii{\sin\bm{\Theta}},\label{eq:geometric-dist}
\end{equation}
where $\vertiii{\cdot}$ is some norm on $\mathbb{R}^{r\times r}$.
\end{enumerate}
\end{enumerate}
As we will soon see, the algebraic distance (\ref{eq:algebraic-dist})
and the geometric distance (\ref{eq:geometric-dist}) are closely
related. This connection is best illustrated via the following lemma.

\begin{lemma}\label{lemma:connection-algebraic-geometric}Denote
\[
k=\begin{cases}
r, & \text{if }2r\leq n,\\
n-r, & \text{if }2r>n.
\end{cases}
\]
Then the singular values of $\bm{U}_{0}\bm{U}_{0}^{\top}-\hat{\bm{U}}_{0}\hat{\bm{U}}_{0}^{\top}$
are given by
\[
\underbrace{\sin\theta_{r},\sin\theta_{r},\sin\theta_{r-1},\sin\theta_{r-1},\cdots,\sin\theta_{r+1-k},\sin\theta_{r+1-k}}_{2k},\;\underbrace{0,0,\cdots,0}_{n-2k}.
\]

Suppose that $2r\leq n$. Then the singular values of $\bm{U}_{0}\bm{U}_{0}^{\top}-\hat{\bm{U}}_{0}\hat{\bm{U}}_{0}^{\top}$
are given by
\[
\underbrace{\sin\theta_{r},\sin\theta_{r},\sin\theta_{r-1},\sin\theta_{r-1},\cdots,\sin\theta_{1},\sin\theta_{1}}_{2r},\;\underbrace{0,0,\cdots,0}_{n-2r}.
\]
Similarly, when $2r>n$, the singular values of $\bm{U}_{0}\bm{U}_{0}^{\top}-\hat{\bm{U}}_{0}\hat{\bm{U}}_{0}^{\top}$
are 
\[
\underbrace{\sin\theta_{r},\sin\theta_{r},\sin\theta_{r-1},\sin\theta_{r-1},\cdots,\sin\theta_{2r+1-n},\sin\theta_{2r+1-n}}_{2(n-r)},\;\underbrace{0,0,\cdots,0}_{2r-n}.
\]
\end{lemma}\begin{proof}Let $\bm{U}_{1}\in\mathbb{R}^{n\times(n-r)}$
(resp.~$\hat{\bm{U}}_{1}\in\mathbb{R}^{n\times(n-r)}$) be an orthonormal
basis for the orthogonal complement to $\bm{U}_{0}$ (reps.~$\hat{\bm{U}}_{0}$).
Since singular values are unitarily invariant, we turn to investigating
the singlur values of 
\[
\left[\begin{array}{c}
\bm{U}_{0}^{\top}\\
\bm{U}_{1}^{\top}
\end{array}\right]\left(\bm{U}_{0}\bm{U}_{0}^{\top}-\hat{\bm{U}}_{0}\hat{\bm{U}}_{0}^{\top}\right)\left[\begin{array}{cc}
\hat{\bm{U}}_{1} & \hat{\bm{U}}_{0}\end{array}\right]=\left[\begin{array}{cc}
\bm{U}_{0}^{\top}\hat{\bm{U}}_{1} & \bm{0}\\
\bm{0} & -\bm{U}_{1}^{\top}\hat{\bm{U}}_{0}
\end{array}\right].
\]
As result, the singular values of $\bm{U}_{0}\bm{U}_{0}^{\top}-\hat{\bm{U}}_{0}\hat{\bm{U}}_{0}^{\top}$
are composed of those of $\bm{U}_{0}^{\top}\hat{\bm{U}}_{1}$ and
$-\bm{U}_{1}^{\top}\hat{\bm{U}}_{0}$. Recall from our discussion
on the relationship between eigenvalues and singular values. We turn
to studying the eigenvalues of $\bm{U}_{0}^{\top}\hat{\bm{U}}_{1}\hat{\bm{U}}_{1}^{\top}\bm{U}_{0}$.
Note that 
\[
\bm{U}_{0}^{\top}\hat{\bm{U}}_{1}\hat{\bm{U}}_{1}^{\top}\bm{U}_{0}=\bm{U}_{0}^{\top}\left(\bm{I}_{n}-\hat{\bm{U}}_{0}\hat{\bm{U}}_{0}^{\top}\right)\bm{U}_{0}=\bm{I}_{r}-\bm{U}_{0}^{\top}\hat{\bm{U}}_{0}\hat{\bm{U}}_{0}^{\top}\bm{U}_{0}=\bm{X}\left(\bm{I}_{r}-\bm{\Sigma}^{2}\right)\bm{X}^{\top},
\]
where the last relation uses the SVD $\bm{X}\bm{\Sigma}\bm{Y}^{\top}$
of $\bm{U}_{0}^{\top}\hat{\bm{U}}_{0}$. It is then easy to check
that 
\[
\sigma_{i}(\bm{U}_{0}^{\top}\hat{\bm{U}}_{1})=\sqrt{\lambda_{i}(\bm{U}_{0}^{\top}\hat{\bm{U}}_{1}\hat{\bm{U}}_{1}^{\top}\bm{U}_{0})}=\sqrt{1-\cos^{2}\theta_{r+1-i}}=\sin\theta_{r+1-i}.
\]
Similaly, the singular values of $-\bm{U}_{1}^{\top}\hat{\bm{U}}_{0}$
are given by $\{\sin\theta_{i}\}_{1\leq i\leq r}$. Combine these
two observations to complete the proof. \end{proof}

Lemma~\ref{lemma:connection-algebraic-geometric} connects the singular
values of the difference of projection matrices with the sines of
canonical angles between subspaces. With the help of Lemma~\ref{lemma:connection-algebraic-geometric},
one can easily obtain the following corollary.

\begin{corollary}\label{prop:unitary-norm-property}One has 
\begin{align*}
\left\Vert \bm{U}_{0}\bm{U}_{0}^{\top}-\hat{\bm{U}}_{0}\hat{\bm{U}}_{0}^{\top}\right\Vert  & =\left\Vert \sin\bm{\Theta}\right\Vert \qquad\text{and}\qquad\left\Vert \bm{U}_{0}\bm{U}_{0}^{\top}-\hat{\bm{U}}_{0}\hat{\bm{U}}_{0}^{\top}\right\Vert _{\mathrm{F}}=\sqrt{2}\left\Vert \sin\bm{\Theta}\right\Vert _{\mathrm{F}}.
\end{align*}
\end{corollary}

\subsubsection{Perturbation of eigenspaces: Davis-Kahan's sin-$\Theta$ theorem}

We now proceed to the perturbation bounds for eigenspaces. Unlike
eigenvalues, which are stable against small perturbations, the eigenspace
of a real symmetric matrix can change drastically even under a small
perturbation. For instance, take
\[
\bm{A}=\left[\begin{array}{cc}
1+\epsilon & 0\\
0 & 1-\epsilon
\end{array}\right],\quad\bm{E}=\left[\begin{array}{cc}
-\epsilon & \epsilon\\
\epsilon & \epsilon
\end{array}\right],\quad\text{and}\quad\hat{\bm{A}}=\bm{A}+\bm{E}=\left[\begin{array}{cc}
1 & \epsilon\\
\epsilon & 1
\end{array}\right],
\]
where $\epsilon$ is a small positive constant. It is straightforward
to compute that the leading eigenvector of $\bm{A}$ (reps.~$\hat{\bm{A}}$)
is $\bm{u}_{1}=(1,0)^{\top}$ (resp.~$\hat{\bm{u}}_{1}=(1/\sqrt{2},1/\sqrt{2})^{\top}$).
As a result, the sin$\bm{\Theta}$ distance between the subspaces
spanned by $\bm{u}_{1}$ and $\hat{\bm{u}}_{1}$ is $1/\sqrt{2}$,
which can be arbitrarily large compared with the small perturbation
$\epsilon$. On closer inspection, this pathological behavior exists
due to the fact that size $\epsilon$ of the perturbation is comparable
to the eigengap of $\bm{A}$, i.e.~$\lambda_{1}(\bm{A})-\lambda_{2}(\bm{A})=2\epsilon$.
This motivates us to impose conditions on the eigengap such that the
eigenspace could be stable against perturbations. The stability of
eigenspaces is formally stated in Theorem~\ref{thm:davis-kahan}.

\begin{theorem}[Davis-Kahan sin-$\Theta$ theorem]\label{thm:davis-kahan}Let
$\bm{A}$ and $\hat{\bm{A}}=\bm{A}+\bm{E}$ be two symmetric matrices
in $\mathbb{R}^{n\times n}$ with the spectral decomposition
\begin{align*}
\bm{A} & =\bm{U}\bm{\Lambda}\bm{U}^{\top}=\left[\begin{array}{cc}
\bm{U}_{0} & \bm{U}_{1}\end{array}\right]\left[\begin{array}{cc}
\bm{\Lambda}_{0} & \bm{0}\\
\bm{0} & \bm{\Lambda}_{1}
\end{array}\right]\left[\begin{array}{c}
\bm{U}_{0}^{\top}\\
\bm{U}_{1}^{\top}
\end{array}\right];\\
\hat{\bm{A}} & =\hat{\bm{U}}\hat{\bm{\Lambda}}\hat{\bm{U}}=\left[\begin{array}{cc}
\hat{\bm{U}}_{0} & \hat{\bm{U}}_{1}\end{array}\right]\left[\begin{array}{cc}
\hat{\bm{\Lambda}}_{0} & \bm{0}\\
\bm{0} & \hat{\bm{\Lambda}}_{1}
\end{array}\right]\left[\begin{array}{c}
\hat{\bm{U}}_{0}^{\top}\\
\hat{\bm{U}}_{1}^{\top}
\end{array}\right],
\end{align*}
respectively. Here 
\begin{align*}
\bm{\Lambda}_{0} & =\mathsf{diag}\left(\lambda_{1},\lambda_{2},\cdots\lambda_{r}\right),\qquad\bm{\Lambda}_{1}=\mathsf{diag}\left(\lambda_{r+1},\lambda_{r+2},\cdots\lambda_{n}\right),\\
\hat{\bm{\Lambda}}_{0} & =\mathsf{diag}\big(\hat{\lambda}_{1},\hat{\lambda}_{2},\cdots\hat{\lambda}_{r}\big),\qquad\hat{\bm{\Lambda}}_{1}=\mathsf{diag}\big(\hat{\lambda}_{r+1},\hat{\lambda}_{r+2},\cdots\hat{\lambda}_{n}\big),
\end{align*}
 and $\lambda_{1}\geq\lambda_{2}\geq\cdots\geq\lambda_{r}>\lambda_{r+1}\geq\cdots\geq\lambda_{n}$,
$\hat{\lambda}_{1}\geq\hat{\lambda}_{2}\geq\cdots\geq\hat{\lambda}_{n}$.
Suppose that the perturbation is not too large in the sense that $\|\bm{E}\|<\lambda_{r}-\lambda_{r+1}$.
Then we have for any unitarily invariant norm $\vertiii{\cdot}$,
\[
\vertiii{\sin\bm{\Theta}\big(\bm{U}_{0},\hat{\bm{U}}_{0}\big)}\leq\frac{\vertiii{\bm{E}\bm{U}_{0}}}{\lambda_{r}-\left(\lambda_{r+1}+\left\Vert \bm{E}\right\Vert \right)}.
\]
\end{theorem}

\begin{theorem}Consider the settings in Section 2.2.1. Assume that
all the eigenvalues of $\bm{\Lambda}^{\star}$ lie in the interval
$[\alpha,\beta]$ and those of $\bm{\Lambda}_{\perp}$ lie outside
of $(\alpha-\Delta,\beta+\Delta)$ for some $\Delta>0$. Then one
has 
\[
\vertiii{\sin\bm{\Theta}}\leq\frac{\vertiii{\bm{E}\bm{U}^{\star}}}{\Delta}.
\]
The same conclusion holds true if the roles of $\bm{\Lambda}^{\star}$
and $\bm{\Lambda}_{\perp}$ are switched. \end{theorem}\begin{proof}Without
loss of generality, we assume that $\alpha=-\beta$ for some $\beta\geq0$.
Two immediate consequences are 
\[
\|\bm{\Lambda}^{\star}\|\leq\beta,\qquad\text{and}\qquad\|\bm{\Lambda}_{\perp}^{-1}\|\leq(\beta+\delta)^{-1}.
\]
Define the residual matrix 
\[
\bm{R}\coloneqq\left(\bm{M}-\bm{M}^{\star}\right)\bm{U}^{\star}=\bm{M}\bm{U}^{\star}-\bm{U}^{\star}\bm{\Lambda}^{\star}.
\]
Multiply both sides by $\bm{U}_{\perp}^{\top}$ to obtain 
\[
\bm{U}_{\perp}^{\top}\bm{R}=\bm{\Lambda}_{\perp}\bm{U}_{\perp}^{\top}\bm{U}^{\star}-\bm{U}_{\perp}^{\top}\bm{U}^{\star}\bm{\Lambda}^{\star}.
\]
Invoke the triangle inequality to see that 
\begin{align*}
\vertiii{\bm{U}_{\perp}^{\top}\bm{R}} & \geq\vertiii{\bm{\Lambda}_{\perp}\bm{U}_{\perp}^{\top}\bm{U}^{\star}}-\vertiii{\bm{U}_{\perp}^{\top}\bm{U}^{\star}\bm{\Lambda}^{\star}}\\
 & \geq\frac{1}{\|\bm{\Lambda}_{\perp}^{-1}\|}\vertiii{\bm{U}_{\perp}^{\top}\bm{U}^{\star}}-\vertiii{\bm{U}_{\perp}^{\top}\bm{U}^{\star}}\|\bm{\Lambda}^{\star}\|\\
 & \geq(\beta+\delta)\vertiii{\bm{U}_{\perp}^{\top}\bm{U}^{\star}}-\beta\vertiii{\bm{U}_{\perp}^{\top}\bm{U}^{\star}}=\delta\vertiii{\bm{U}_{\perp}^{\top}\bm{U}^{\star}}.
\end{align*}
Here the middle line uses the facts that $\vertiii{\bm{U}_{\perp}^{\top}\bm{U}^{\star}}=\vertiii{\bm{\Lambda}_{\perp}^{-1}\bm{\Lambda}_{\perp}\bm{U}_{\perp}^{\top}\bm{U}^{\star}}\leq\|\bm{\Lambda}_{\perp}^{-1}\|\vertiii{\bm{\Lambda}_{\perp}\bm{U}_{\perp}^{\top}\bm{U}^{\star}}$
and that $\vertiii{\bm{U}_{\perp}^{\top}\bm{U}^{\star}\bm{\Lambda}^{\star}}\leq\vertiii{\bm{U}_{\perp}^{\top}\bm{U}^{\star}}\|\bm{\Lambda}^{\star}\|$;
see Prop XXX. The last inequality arises from the estimates (XXX).
As a consequence, we arrive at the conclusion that 
\[
\vertiii{\bm{U}_{\perp}^{\top}\bm{U}^{\star}}\leq\frac{\vertiii{\bm{U}_{\perp}^{\top}\bm{R}}}{\delta}\leq\frac{\vertiii{\bm{R}}}{\delta}=\frac{\vertiii{\bm{E}\bm{U}^{\star}}}{\delta}.
\]
Identifying $\vertiii{\bm{U}_{\perp}^{\top}\bm{U}^{\star}}$ with
$\vertiii{\sin\bm{\Theta}}$ concludes the proof. \end{proof}

\begin{corollary}Consider the setting where $\bm{\Lambda}^{\star}$
contains all nonzero eigenvalues of $\bm{M}^{\star}$. Accordingly,
assume that $\bm{\Lambda}$ contains the $k$ largest eigenvalues
of $\bm{M}$ in the absolute value sense. Then if $\|\bm{E}\|\leq1/2\min_{1\leq i\leq k}|\lambda_{i}^{\star}|$,
one has 
\[
\vertiii{\sin\bm{\Theta}\big(\bm{U}^{\star},\bm{U}\big)}\leq\frac{\vertiii{\bm{E}\bm{U}^{\star}}}{\min_{1\leq i\leq k}|\lambda_{i}^{\star}|-\|\bm{E}\|}.
\]
 \end{corollary}\begin{proof}In view of Weyl's inequality, we know
that for $i=1,2,\cdots,n$
\[
|\lambda_{i}(\bm{M})-\lambda_{i}(\bm{M}^{\star})|\leq\|\bm{E}\|\leq\frac{1}{2}\min_{1\leq i\leq k}|\lambda_{i}^{\star}|,
\]
which immediately implies that all the eigenvalues of $\bm{\Lambda}_{\perp}$
lie within the interval $[-\|\bm{E}\|,\|\bm{E}\|]$. This together
with the assumption that all the eigenvalues of $\bm{\Lambda}^{\star}$
are in $\mathbb{R}\backslash(-\min_{1\leq i\leq k}|\lambda_{i}^{\star}|,\min_{1\leq i\leq k}|\lambda_{i}^{\star}|)$
allows us to invoke Theorem XXX to obtain
\[
\vertiii{\sin\bm{\Theta}}\leq\frac{\vertiii{\bm{E}\bm{U}^{\star}}}{\min_{1\leq i\leq k}|\lambda_{i}^{\star}|-\|\bm{E}\|}.
\]
\end{proof}

\begin{proof}In view of the decomposition $\bm{A}=\bm{U}_{0}\bm{\Lambda}_{0}\bm{U}_{0}^{\top}+\bm{U}_{1}\bm{\Lambda}_{1}\bm{U}_{1}^{\top}$
and the fact that $\bm{U}_{1}^{\top}\bm{U}_{0}=\bm{0}$, we have 
\[
\bm{A}\bm{U}_{0}=\bm{U}_{0}\bm{\Lambda}_{0},
\]
which further implies 
\begin{equation}
\hat{\bm{U}}_{1}^{\top}\bm{A}\bm{U}_{0}=\hat{\bm{U}}_{1}^{\top}\bm{U}_{0}\bm{\Lambda}_{0}.\label{eq:dk-identity-1}
\end{equation}
In addition, the similar decomposition $\bm{A}+\bm{E}=\hat{\bm{U}}_{0}\hat{\bm{\Lambda}}_{0}\hat{\bm{U}}_{0}^{\top}+\hat{\bm{U}}_{1}\hat{\bm{\Lambda}}_{1}\hat{\bm{U}}_{1}^{\top}$
and the fact that $\hat{\bm{U}}_{0}^{\top}\hat{\bm{U}}_{1}=\bm{0}$
tell us that $(\bm{A}+\bm{E})\hat{\bm{U}}_{1}=\hat{\bm{U}}_{1}\hat{\bm{\Lambda}}_{1}$
and hence 
\begin{equation}
\bm{U}_{0}^{\top}\left(\bm{A}+\bm{E}\right)\hat{\bm{U}}_{1}=\bm{U}_{0}^{\top}\hat{\bm{U}}_{1}\hat{\bm{\Lambda}}_{1}.\label{eq:dk-identity-2}
\end{equation}
Combine the two identities (\ref{eq:dk-identity-1}) and (\ref{eq:dk-identity-2})
to reach 
\begin{align}
\hat{\bm{U}}_{1}^{\top}\bm{E}\bm{U}_{0} & =\hat{\bm{\Lambda}}_{1}\hat{\bm{U}}_{1}^{\top}\bm{U}_{0}-\hat{\bm{U}}_{1}^{\top}\bm{U}_{0}\bm{\Lambda}_{0}\nonumber \\
 & =\big(\hat{\bm{\Lambda}}_{1}-c\bm{I}_{n-r}\big)\hat{\bm{U}}_{1}^{\top}\bm{U}_{0}-\hat{\bm{U}}_{1}^{\top}\bm{U}_{0}\left(\bm{\Lambda}_{0}-c\bm{I}_{r}\right),\label{eq:dk-identity-main}
\end{align}
where the last line holds for any $c\in\mathbb{R}$. In particular,
we can take $c\triangleq(\lambda_{1}+\lambda_{r})/2$ and denote $\rho\triangleq(\lambda_{1}-\lambda_{r})/2\geq0$.
As a result, all the eigenvalues of $\bm{\Lambda}_{0}-c\bm{I}_{r}$
are contained in the range $[-\rho,\rho]$ whereas all the eigenvalues
of $\hat{\bm{\Lambda}}_{1}-c\bm{I}_{n-r}$ are smaller than $\hat{\lambda}_{r+1}-c$.
Apply Weyl's inequality (see Lemma~\ref{lemma:weyl}) to see that
$\hat{\lambda}_{r+1}\leq\lambda_{r+1}+\|\bm{E}\|$. This combined
with relation $\lambda_{r+1}-c=\lambda_{r+1}-\lambda_{r}-\rho$ reveals
that all the eigenvalues of $\hat{\bm{\Lambda}}_{1}-c\bm{I}_{n-r}$
are within $(-\infty,\lambda_{r+1}-\lambda_{r}+\|\bm{E}\|-\rho)$.
Two immediate consequences are that 
\begin{equation}
\left\Vert \bm{\Lambda}_{0}-c\bm{I}_{r}\right\Vert \le\rho,\qquad\text{and}\qquad\sigma_{\min}\big(\hat{\bm{\Lambda}}_{1}-c\bm{I}_{n-r}\big)\geq\rho+\lambda_{r}-\lambda_{r+1}-\left\Vert \bm{E}\right\Vert .\label{eq:dk-relation-main}
\end{equation}
Continue the derivation in (\ref{eq:dk-identity-main}) to obtain
\begin{align*}
\vertiii{\hat{\bm{U}}_{1}^{\top}\bm{E}\bm{U}_{0}} & \overset{(\text{i})}{\geq}\vertiii{\left(\hat{\bm{\Lambda}}_{1}-c\bm{I}_{n-r}\right)\hat{\bm{U}}_{1}^{\top}\bm{U}_{0}}-\vertiii{\hat{\bm{U}}_{1}^{\top}\bm{U}_{0}\left(\bm{\Lambda}_{0}-c\bm{I}_{r}\right)}\\
 & \overset{(\text{ii})}{\geq}\left(\rho+\lambda_{r}-\lambda_{r+1}-\left\Vert \bm{E}\right\Vert _{\mathsf{}}\right)\vertiii{\hat{\bm{U}}_{1}^{\top}\bm{U}_{0}}-\rho\vertiii{\hat{\bm{U}}_{1}^{\top}\bm{U}_{0}}\\
 & =\left(\lambda_{r}-\lambda_{r+1}-\left\Vert \bm{E}\right\Vert _{\mathsf{}}\right)\vertiii{\hat{\bm{U}}_{1}^{\top}\bm{U}_{0}},
\end{align*}
where (i) follows from the triangle inequality and (ii) holds because
of Proposition~\ref{prop:unitary-norm-property} and the conditions
(\ref{eq:dk-relation-main}). Clearly, the above inequality is equivalent
to 
\[
\vertiii{\hat{\bm{U}}_{1}^{\top}\bm{U}_{0}}\leq\frac{\vertiii{\hat{\bm{U}}_{1}^{\top}\bm{E}\bm{U}_{0}}}{\lambda_{r}-\lambda_{r+1}-\left\Vert \bm{E}\right\Vert }\leq\frac{\vertiii{\bm{E}\bm{U}_{0}}}{\lambda_{r}-\left(\lambda_{r+1}+\left\Vert \bm{E}\right\Vert \right)}.
\]
The proof is then completed via observing Lemma~\ref{lemma:sin-theta-equivalence}.\end{proof}

\begin{lemma}\label{lemma:sin-theta-equivalence}Suppose that $[\bm{U}_{0},\bm{U}_{1}]$
and $[\hat{\bm{U}}_{0},\hat{\bm{U}}_{1}]$ are orthonormal matrices.
Let $\{\theta_{i}\}_{1\leq i\leq r}$ be the canonical angles between
the space spanned by $\bm{U}_{0}$ and that by $\hat{\bm{U}}_{0}$.
Then the singular values of $\hat{\bm{U}}_{1}^{\top}\bm{U}_{0}$ are
given by 
\[
\sin\theta_{r},\sin\theta_{r-1},\cdots\sin\theta_{1},\;0,0,\cdots0.
\]
Consequently, one has 
\begin{equation}
\left\Vert \sin\bm{\Theta}\left(\bm{U}_{0},\hat{\bm{U}}_{0}\right)\right\Vert =\left\Vert \hat{\bm{U}}_{1}^{\top}\bm{U}_{0}\right\Vert \qquad\text{and}\qquad\left\Vert \sin\bm{\Theta}\left(\bm{U}_{0},\hat{\bm{U}}_{0}\right)\right\Vert _{\mathrm{F}}=\left\Vert \hat{\bm{U}}_{1}^{\top}\bm{U}_{0}\right\Vert _{\mathrm{F}}.\label{eq:sin-theta-equivalence}
\end{equation}
Moreover, for any unitarily invariant norm $\vertiii{\cdot}$, we
have 
\[
\vertiii{\sin\bm{\Theta}\left(\bm{U}_{0},\hat{\bm{U}}_{0}\right)}=\vertiii{\hat{\bm{U}}_{1}^{\top}\bm{U}_{0}}.
\]
\end{lemma}\begin{proof}Consider the matrix 
\[
\bm{U}_{0}^{\top}\hat{\bm{U}}_{1}\hat{\bm{U}}_{1}^{\top}\bm{U}_{0}=\bm{U}_{0}^{\top}\big(\bm{I}_{n}-\hat{\bm{U}}_{0}\hat{\bm{U}}_{0}^{\top}\big)\bm{U}_{0}=\bm{I}_{r}-\bm{U}_{0}^{\top}\hat{\bm{U}}_{0}\hat{\bm{U}}_{0}^{\top}\bm{U}_{0}.
\]
Let $\bm{X}\bm{\Sigma}\bm{Y}^{\top}$ be the SVD of $\bm{U}_{0}^{\top}\hat{\bm{U}}_{0}$.
By definition, one has $\bm{\Sigma}=\mathsf{diag}(\cos\theta_{1},\cos\theta_{2},\cdots,\cos\theta_{r})$.
Continue the derivation above to obtain 
\[
\bm{U}_{0}^{\top}\hat{\bm{U}}_{1}\hat{\bm{U}}_{1}^{\top}\bm{U}_{0}=\bm{I}_{r}-\bm{X}\bm{\Sigma}^{2}\bm{X}^{\top}=\bm{X}\left(\bm{I}_{r}-\bm{\Sigma}^{2}\right)\bm{X}^{\top}.
\]
Utilizing the relationship between singular values and eigenvalues,
one sees that for $1\leq i\leq r$
\[
\sigma_{i}\big(\hat{\bm{U}}_{1}^{\top}\bm{U}_{0}\big)=\sqrt{\lambda_{i}\left(\bm{I}_{r}-\bm{\Sigma}^{2}\right)}=\sqrt{\lambda_{i}\left(1-\cos^{2}\theta_{r-i+1}\right)}=\sin\theta_{r-i+1}.
\]
This finishes the first part. The consequences in (\ref{eq:sin-theta-equivalence})
follow immediately. The last relation needs more care. XXX \cm{TODO}\end{proof}

\subsection{Perturbation of general matrices}

Moving beyond real symmetric matrices, we have similar perturbation
bounds on singular values and singular subspaces, parallel to Lemma~\ref{lemma:weyl}
and Theorem~\ref{thm:davis-kahan}.

\subsubsection{Perturbation of singular values}

We begin with the perturbation bounds for singular values.

\begin{lemma}[Wely's inequality for singular values]\label{lemma:weyls-singular-value}Let
$\bm{A},\bm{E}\in\mathbb{R}^{m\times n}$ be two general matrices.
Without loss of generality, assume that $m\leq n$. Then for every
$1\leq i\leq m$, one has 
\[
\left|\sigma_{i}\left(\bm{A}+\bm{E}\right)-\sigma_{i}\left(\bm{A}\right)\right|\leq\left\Vert \bm{E}\right\Vert .
\]
\end{lemma}

Again, Lemma~\ref{lemma:weyls-singular-value} reveals that singular
values are \emph{stable }with respect to small perturbations.

\subsubsection{Perturbation of singular subspaces: Wedin's sin-$\Theta$ theorem}

One has similar perturbation bounds for singular subspaces, as proposed
by Wedin. \cm{Add citations.}

\begin{theorem}[Wedin's sin-$\Theta$ theorem]\label{thm:wedin}Let
$\bm{A}$ and $\hat{\bm{A}}=\bm{A}+\bm{E}$ be two matrices in $\mathbb{R}^{m\times n}$
$(m\leq n)$ with the singular value decomposition
\begin{align*}
\bm{A} & =\bm{U}\bm{\Sigma}\bm{V}^{\top}=\left[\begin{array}{cc}
\bm{U}_{0} & \bm{U}_{1}\end{array}\right]\left[\begin{array}{ccc}
\bm{\Sigma}_{0} & \bm{0} & \bm{0}\\
\bm{0} & \bm{\Sigma}_{1} & \bm{0}
\end{array}\right]\left[\begin{array}{c}
\bm{V}_{0}^{\top}\\
\bm{V}_{1}^{\top}
\end{array}\right];\\
\hat{\bm{A}} & =\hat{\bm{U}}\hat{\bm{\Sigma}}\hat{\bm{V}}^{\top}=\left[\begin{array}{cc}
\hat{\bm{U}}_{0} & \hat{\bm{U}}_{1}\end{array}\right]\left[\begin{array}{ccc}
\hat{\bm{\Sigma}}_{0} & \bm{0} & \bm{0}\\
\bm{0} & \hat{\bm{\Sigma}}_{1} & \bm{0}
\end{array}\right]\left[\begin{array}{c}
\hat{\bm{V}}_{0}^{\top}\\
\hat{\bm{V}}_{1}^{\top}
\end{array}\right],
\end{align*}
respectively. Here 
\begin{align*}
\bm{\Sigma}_{0} & =\mathsf{diag}\left(\sigma_{1}\left(\bm{A}\right),\sigma_{2}\left(\bm{A}\right),\cdots,\sigma_{r}\left(\bm{A}\right)\right),\qquad\bm{\Sigma}_{1}=\mathsf{diag}\left(\sigma_{r+1}\left(\bm{A}\right),\sigma_{r+2}\left(\bm{A}\right),\cdots,\sigma_{n}\left(\bm{A}\right)\right),\\
\hat{\bm{\Sigma}}_{0} & =\mathsf{diag}\big(\sigma_{1}\big(\hat{\bm{A}}\big),\sigma_{2}\big(\hat{\bm{A}}\big),\cdots,\sigma_{r}\big(\hat{\bm{A}}\big)\big),\qquad\hat{\bm{\Sigma}}_{1}=\mathsf{diag}\big(\sigma_{r+1}\big(\hat{\bm{A}}\big),\sigma_{r+2}\big(\hat{\bm{A}}\big),\cdots,\sigma_{n}\big(\hat{\bm{A}}\big)\big).
\end{align*}
 Suppose that the perturbation is not too large in the sense that
$\|\bm{E}\|<\sigma_{r}(\bm{A})-\sigma_{r+1}(\bm{A})$. Then one has
for any unitarily invariant norm $\vertiii{\cdot}$, 
\[
\max\left\{ \vertiii{\sin\bm{\Theta}\left(\bm{U}_{0},\hat{\bm{U}}_{0}\right)},\vertiii{\sin\bm{\Theta}\left(\bm{V}_{0},\hat{\bm{V}}_{0}\right)}\right\} \leq\frac{\max\left\{ \vertiii{\bm{E}^{\top}\bm{U}_{0}},\vertiii{\bm{E}\bm{V}_{0}}\right\} }{\sigma_{r}\left(\bm{A}\right)-\left(\sigma_{r+1}\left(\bm{A}\right)+\|\bm{E}\|\right)}.
\]
\end{theorem}\begin{proof}\cm{Leave it to the reader since its
very similar to the proof of the Davis-Kahan theorem.}\end{proof}

\section{Notes}

In this chapter, we have reviewed classical matrix perturbation theory
that are useful for our subsequent developments on spectral methods.
We conclude this chapter with additional references on this subject. 

Our exposition largely follows the excellent book by Stewart \cite{stewart1990matrix}.
Two other classical references are \cite{bhatia2013matrix,horn2012matrix},
which go beyond matrix perturbation theory and deal with general matrix
analysis. More broadly, matrices can be viewed as linear operators
on the finite dimensional Euclidean space; the general treatment of
perturbation theory for linear operators can be found in \cite{kato2013perturbation}. 

Recall that in Lemma 2.2 \cm{Need to change later} we have stated
the equivalence between $\|\bm{U}_{\perp}^{\top}\bm{U}^{\star}\|$
(resp.~$\|\bm{U}_{\perp}^{\top}\bm{U}^{\star}\|_{\mathrm{F}}$) and
$\|\sin\bm{\Theta}\|$ (resp.~$\|\sin\bm{\Theta}\|_{\mathrm{F}}$).
It turns out that $\vertiii{\bm{U}_{\perp}^{\top}\bm{U}^{\star}}=\vertiii{\sin\bm{\Theta}}$
holds for any unitarily invariant norm; see \cite[Lemma 2.2]{li1998relative}.
In addition, \cite[Theorems 3.2 and 3.4]{li1998relative} offer general
versions of Davis-Kahan's and Wedin's theorems, which do not restrict
attention to the top eigenspaces or singular subspaces. A user-friendly
version of the Davis-Kahan theorem was provided in \cite{MR3371006}. 

Consider the Procrustes problem 
\[
\mathsf{minimize}_{\bm{R}\in\mathcal{O}^{r\times r}}\quad\left\Vert \hat{\bm{U}}\bm{R}-\bm{U}\right\Vert _{\mathrm{F}}^{2},
\]
whose solution is revealed in the SVD of $\hat{\bm{U}}^{\top}\bm{U}$.
More specifically, let $\bm{X}\bm{\Sigma}\bm{Y}^{\top}$ be the SVD
of $\hat{\bm{U}}^{\top}\bm{U}$, then one has 
\[
\bm{X}\bm{Y}^{\top}=\arg\min_{\bm{R}\in\mathcal{O}^{r\times r}}\quad\left\Vert \hat{\bm{U}}\bm{R}-\bm{U}\right\Vert _{\mathrm{F}}^{2}.
\]
Denote by $(\sigma_{1},\sigma_{2},\cdots\sigma_{r})$ the entries
in $\bm{\Sigma}$ in decreasing order. It is easy to check that they
are contained in $[0,1]$. As a result, we define the principal angles
between two subspaces spaned by $\bm{U}$ and $\hat{\bm{U}}$ to be
\[
\theta_{i}\triangleq\cos^{-1}\left(\sigma_{i}\right),\qquad\text{for all}\qquad1\leq i\leq r
\]
and arrange them into a diagonal matrix 
\[
\bm{\Theta}\triangleq\mathsf{diag}\left(\theta_{1},\theta_{2},\cdots,\theta_{r}\right)\in\mathbb{R}^{r\times r}.
\]
With $\bm{\Theta}$ in place, we can define 
\begin{align*}
\cos\bm{\Theta} & =\mathsf{diag}\left(\cos\theta_{1},\cos\theta_{2},\cdots,\cos\theta_{r}\right)=\bm{\Sigma};\\
\sin\bm{\Theta} & =\mathsf{diag}\left(\sin\theta_{1},\sin\theta_{2},\cdots,\sin\theta_{r}\right).
\end{align*}

Several immediate facts are in order. First 
\[
\left\Vert \sin\bm{\Theta}\right\Vert _{\mathrm{F}}^{2}=\sum_{i=1}^{r}\left(\sin\theta_{r}\right)^{2}=\sum_{i=1}^{r}\left[1-\left(\cos\theta_{r}\right)^{2}\right]=r-\left\Vert \bm{\Sigma}\right\Vert _{\mathrm{F}}^{2}.
\]
Second, one has 
\begin{align*}
\left\Vert \hat{\bm{U}}\bm{X}\bm{Y}^{\top}-\bm{U}\right\Vert _{\mathrm{F}}^{2} & =\left\Vert \hat{\bm{U}}\right\Vert _{\mathrm{F}}^{2}+\left\Vert \bm{U}\right\Vert _{\mathrm{F}}^{2}-2\mathsf{Tr}\left(\bm{Y}\bm{X}^{\top}\hat{\bm{U}}^{\top}\bm{U}\right)\\
 & =2r-2\mathsf{Tr}\left(\bm{Y}\bm{X}^{\top}\bm{X}\bm{\Sigma}\bm{Y}^{\top}\right)\\
 & =2r-2\mathsf{Tr}\left(\bm{\Sigma}\right)\\
 & \leq2r-2\sum_{i=1}^{r}\sigma_{i}^{2}\\
 & =2\left\Vert \sin\bm{\Theta}\right\Vert _{\mathrm{F}}^{2}.
\end{align*}
Last but not least, the distance between projection matrices satisfy
\begin{align*}
\left\Vert \hat{\bm{U}}\hat{\bm{U}}^{\top}-\bm{U}\bm{U}^{\top}\right\Vert _{\mathrm{F}}^{2} & =\left\Vert \hat{\bm{U}}\hat{\bm{U}}^{\top}\right\Vert _{\mathrm{F}}^{2}+\left\Vert \hat{\bm{U}}\hat{\bm{U}}^{\top}\right\Vert _{\mathrm{F}}^{2}-2\mathsf{Tr}\left(\hat{\bm{U}}\hat{\bm{U}}^{\top}\bm{U}\bm{U}^{\top}\right)\\
 & =2r-2\left\Vert \hat{\bm{U}}^{\top}\bm{U}\right\Vert _{\mathrm{F}}^{2}\\
 & =2r-2\left\Vert \bm{\Sigma}\right\Vert _{\mathrm{F}}^{2}\\
 & =2\left\Vert \sin\bm{\Theta}\right\Vert _{\mathrm{F}}^{2}.
\end{align*}


\section{Matrix completion}

Imagine that one observes a small subset of the entries of a matrix
and seeks to recover the full matrix. \cm{Add a few words on applications.}

\subsection{Setup}

Let $\bm{M}^{\star}\in\mathbb{R}^{n_{1}\times n_{2}}$ be the rank-$r$
matrix we hope to recover. Without loss of generality, we assume that
$n_{1}\leq n_{2}$. Denote by $\Omega\subseteq\{1,2\cdots n_{1}\}]\times\{1,2,\cdots,n_{2}\}$
the index set for the observed entries, and correpondingly define
the projection operator $\mathcal{P}_{\Omega}:\mathbb{R}^{n_{1}\times n_{2}}\mapsto\mathbb{R}^{n_{1}\times n_{2}}$
to be 
\[
[\mathcal{P}_{\Omega}(\bm{A})]_{i,j}=\begin{cases}
A_{i,j}, & \text{if }(i,j)\in\Omega,\\
0, & \text{else},
\end{cases}
\]
where $\bm{A}\in\mathbb{R}^{n_{1}\times n_{2}}$ denotes the input
matrix. With this notation in place, matrix completion amounts to
recovering $\bm{M}^{\star}$ from its partial entries $\mathcal{P}_{\Omega}(\bm{M}^{\star})$.
Throughout this paper, we impose the Bernoulli observation model detailed
in the following assumption. 

\begin{assumption}[Bernoulli observation model for matrix completion]We
assume that each entry in $\bm{M}^{\star}$ is observed independentlywith
probability $0<p\leq1$, i.e.~$(i,j)\in\Omega$ independently with
probability $p$ for each $1\leq i\leq n_{1}$ and $1\leq j\leq n_{2}$.
\end{assumption}

The random observation model alone does not guarantee effective recovery
of the underlying matrix $\bm{M}^{\star}$. It turns out the following
incoherence property of $\bm{M}^{\star}$ plays a crutial role in
recovering low-rank matrices from partial observations. 

\begin{definition}[Incoherence for matrix completion]A rank-$r$
matrix $\bm{M}^{\star}\in\mathbb{R}^{n_{1}\times n_{2}}$ with compact
SVD $\bm{M}^{\star}=\bm{U}^{\star}\bm{\Sigma}^{\star}\bm{V}^{\star}$
is said to be $\mu$-incoherent for some $\mu\geq1$ if 
\[
\|\bm{U}^{\star}\|_{2,\infty}\leq\sqrt{\frac{\mu}{n_{1}}}\|\bm{U}^{\star}\|_{\mathrm{F}}=\sqrt{\frac{\mu r}{n_{1}}},\quad\text{and}\quad\|\bm{V}^{\star}\|_{2,\infty}\leq\sqrt{\frac{\mu}{n_{2}}}\|\bm{V}^{\star}\|_{\mathrm{F}}=\sqrt{\frac{\mu r}{n_{2}}}.
\]
\end{definition}

Before introducing the spectral method for matrix completion, we single
out a few immediate consequences of $\mu$-incoherence. First, the
maximum row norm $\|\bm{M}^{\star}\|_{2,\infty}$ of $\bm{M}^{\star}$
cannot be two large, i.e. 
\begin{equation}
\|\bm{M}^{\star}\|_{2,\infty}=\|\bm{U}^{\star}\bm{\Sigma}^{\star}\bm{V}^{\star\top}\|_{2,\infty}\leq\|\bm{U}^{\star}\|_{2,\infty}\|\bm{\Sigma}^{\star}\|\|\bm{V}^{\star}\|\leq\sqrt{\frac{\mu r}{n_{1}}}\sigma_{1}(\bm{M}^{\star}),\label{eq:mc-row-norm-upper}
\end{equation}
where the first inequality arises from the elementary inequalities
$\|\bm{A}\bm{B}\|_{2,\infty}\leq\|\bm{A}\|_{2,\infty}\|\bm{B}\|$
and $\|\bm{A}\bm{B}\|\leq\|\bm{A}\|\|\bm{B}\|$. The last relation
uses the incoherence assumption on $\bm{M}^{\star}$, the orthogonality
of $\bm{V}^{\star}$, and identifies $\|\bm{\Sigma}^{\star}\|$ with
$\sigma_{1}(\bm{M}^{\star})$. An analogous upper bound holds for
$\|\bm{M}^{\star\top}\|_{2,\infty}$:
\begin{equation}
\|\bm{M}^{\star\top}\|_{2,\infty}\|\bm{V}^{\star}\bm{\Sigma}^{\star}\bm{U}^{\star\top}\|_{2,\infty}\leq\|\bm{V}^{\star}\|_{2,\infty}\|\bm{\Sigma}^{\star}\|\|\bm{U}^{\star}\|\leq\sqrt{\frac{\mu r}{n_{2}}}\sigma_{1}(\bm{M}^{\star}).\label{eq:mc-column-norm-upper}
\end{equation}
In addition, the matrix $\bm{M}^{\star}$ is element-wise bounded
in the sense that 
\begin{equation}
\|\bm{M}^{\star}\|_{\infty}=\|\bm{U}^{\star}\bm{\Sigma}^{\star}\bm{V}^{\star\top}\|_{\infty}\leq\|\bm{U}^{\star}\|_{2,\infty}\|\bm{\Sigma}^{\star}\|\|\bm{V}^{\star}\|_{2,\infty}\leq\frac{\mu r}{\sqrt{n_{1}n_{2}}}\sigma_{1}(\bm{M}^{\star}).\label{eq:mc-entry-upper}
\end{equation}
Here, the first inequality follows from the relation $\|\bm{A}\bm{B}^{\top}\|_{\infty}\leq\|\bm{A}\|_{2,\infty}\|\bm{B}\|_{2,\infty}$
and the aforementioned $\|\bm{A}\bm{B}\|_{2,\infty}\leq\|\bm{A}\|_{2,\infty}\|\bm{B}\|$,
while the last one again utilizes the incoherence assumption. 

\subsection{Spectral method for matrix completion}

Recall that the idea of spectral method is to first form a good approximation
$\bm{M}$ to the unknown matrix $\bm{M}^{\star}$ we care about, and
then perform spectral decomposition on $\bm{M}$ to approximately
recover the singular spaces (spectral properties) of $\bm{M}^{\star}$.
Given our random sampling model, such a good approximation is readily
available: 
\[
\bm{M}\coloneqq p^{-1}\mathcal{P}_{\Omega}(\bm{M}^{\star}).
\]
Indeed, one can check that $\mathbb{E}[\bm{M}]=\bm{M}^{\star}$ where
the expectation is taken over the randomness in $\Omega$. As a result,
it is reasonble to use $\bm{U}\bm{\Sigma}\bm{V}^{\top}$, the top-$r$
SVD of $p^{-1}\mathcal{P}_{\Omega}(\bm{M}^{\star})$ as an approximation
to the groundtruth $\bm{M}^{\star}$. We learned from Wedin's $\sin\Theta$
theorem that whether or not $\bm{U}$ (resp.~$\bm{V}$) is close
to $\bm{U}^{\star}$ (resp.~$\bm{V}^{\star}$) relies crutially on
the size of the perturbation $\|p^{-1}\mathcal{P}_{\Omega}(\bm{M}^{\star})-\bm{M}^{\star}\|$.
Therefore we begin with offering an upper bound on $\|p^{-1}\mathcal{P}_{\Omega}(\bm{M}^{\star})-\bm{M}^{\star}\|$. 

\begin{lemma}Suppose that $n_{2}p\geq C\mu r\log n_{2}$ for some
sufficiently large constant $C>0$. Then with probability at least
$1-O(n_{2}^{-10})$, one has 
\[
\|p^{-1}\mathcal{P}_{\Omega}(\bm{M}^{\star})-\bm{M}^{\star}\|\lesssim\sqrt{\frac{\mu r}{n_{1}p}\log n_{2}}\,\sigma_{1}(\bm{M}^{\star}).
\]
\end{lemma}\begin{proof}Note that the term $p^{-1}\mathcal{P}_{\Omega}(\bm{M}^{\star})-\bm{M}^{\star}$
can be expressed as the sum of $n_{1}n_{2}$ i.i.d.~random matrices:
\[
p^{-1}\mathcal{P}_{\Omega}(\bm{M}^{\star})-\bm{M}^{\star}=\sum_{i=1}^{n_{1}}\sum_{j=1}^{n_{2}}\underbrace{(p^{-1}\mathbbm{1}\{\delta_{ij}=1\}-1)M_{ij}^{\star}\bm{e}_{i}\bm{e}_{j}^{\top}}_{\eqqcolon\bm{X}_{ij}},
\]
where $\delta_{ij}$ follows an independent Bernoulli distribution
with parameter $p$, and $\bm{e}_{i}$ stands for the $i$-th standard
basis vector. It is easily seen that for each $i,j$ 
\[
\mathbb{E}[\bm{X}_{ij}]=\bm{0},\quad\text{and}\quad\|\bm{X}_{ij}\|\leq p^{-1}\|\bm{M}^{\star}\|_{\infty}\leq\frac{\mu r}{p\sqrt{n_{1}n_{2}}}\sigma_{1}(\bm{M}^{\star}),
\]
where the last relation uses the entrywise upper bound (\ref{eq:mc-entry-upper})
on $\bm{M}^{\star}$. In order to apply the matrix Bernstein inequality,
we need to control the variance statistics
\[
v\coloneqq\max\left\{ \left\Vert \sum_{i=1}^{n_{1}}\sum_{j=1}^{n_{2}}\mathbb{E}\left[\bm{X}_{ij}\bm{X}_{ij}^{\top}\right]\right\Vert ,\left\Vert \sum_{i=1}^{n_{1}}\sum_{j=1}^{n_{2}}\mathbb{E}\left[\bm{X}_{ij}^{\top}\bm{X}_{ij}\right]\right\Vert \right\} .
\]
Regarding the first variance term, we have 
\begin{align*}
\sum_{i=1}^{n_{1}}\sum_{j=1}^{n_{2}}\mathbb{E}\left[\bm{X}_{ij}\bm{X}_{ij}^{\top}\right] & =\sum_{i=1}^{n_{1}}\sum_{j=1}^{n_{2}}\mathbb{E}\left[(p^{-1}\mathbbm{1}\{\delta_{ij}=1\}-1)^{2}(M_{ij}^{\star})^{2}\bm{e}_{i}\bm{e}_{j}^{\top}\bm{e}_{j}\bm{e}_{i}^{\top}\right]\\
 & =\frac{1-p}{p}\sum_{i=1}^{n_{1}}\sum_{j=1}^{n_{2}}(M_{ij}^{\star})^{2}\bm{e}_{i}\bm{e}_{i}^{\top}=\frac{1-p}{p}\sum_{i=1}^{n_{1}}\|\bm{M}_{i,\cdot}^{\star}\|_{2}^{2}\bm{e}_{i}\bm{e}_{i}^{\top}.
\end{align*}
Here, the first identity arises from the definition of $\bm{X}_{ij}$,
the second one calculated the variance of a Bernoulli random variable,
and the last equation follows from a direct calculation. Similarly,
the second term in the variance statistics enjoys the following characterization:
\[
\sum_{i=1}^{n_{1}}\sum_{j=1}^{n_{2}}\mathbb{E}\left[\bm{X}_{ij}^{\top}\bm{X}_{ij}\right]=\frac{1-p}{p}\sum_{j=1}^{n_{2}}\|\bm{M}_{\cdot,j}^{\star}\|_{2}^{2}\bm{e}_{j}\bm{e}_{j}^{\top}.
\]
Taking the above two relations together, we obtain the following upper
bound 
\begin{align*}
v & \leq\frac{1}{p}\max\left\{ \|\bm{M}^{\star}\|_{2,\infty}^{2},\|\bm{M}^{\star\top}\|_{2,\infty}^{2}\right\} \leq\frac{1}{p}\left\{ \frac{\mu r}{n_{1}}\sigma_{1}^{2}(\bm{M}^{\star}),\frac{\mu r}{n_{2}}\sigma_{1}^{2}(\bm{M}^{\star})\right\} \leq\frac{\mu r}{n_{1}p}\sigma_{1}^{2}(\bm{M}^{\star}),
\end{align*}
where the middle relation stems from the upper bounds (\ref{eq:mc-row-norm-upper})
and (\ref{eq:mc-column-norm-upper}) and the last inequality uses
the assumption $n_{1}\leq n_{2}$. Now we are in position to invoke
the matrix Bernstein inequality and arrive at the conclusion that
\[
\|p^{-1}\mathcal{P}_{\Omega}(\bm{M}^{\star})-\bm{M}^{\star}\|\lesssim\sqrt{\frac{\mu r}{n_{1}p}\sigma_{1}^{2}(\bm{M}^{\star})\log n_{2}}+\frac{\mu r}{p\sqrt{n_{1}n_{2}}}\sigma_{1}(\bm{M}^{\star})\log n_{2}
\]
holds with probability at least $1-O(n^{-10})$. To further simplify
the matter, note that when $n_{2}p\geq C\mu r\log n_{2}$ for some
sufficiently large constant $C>0$, one has 
\[
\frac{\mu r}{p\sqrt{n_{1}n_{2}}}\sigma_{1}(\bm{M}^{\star})\log n_{2}\leq\sqrt{\frac{\mu r}{n_{1}p}\sigma_{1}^{2}(\bm{M}^{\star})\log n_{2}},
\]
which further implies the desired claim. \end{proof}

As a consequence of Lemma XXX, one has 
\[
\|p^{-1}\mathcal{P}_{\Omega}(\bm{M}^{\star})-\bm{M}^{\star}\|\lesssim\sqrt{\frac{\mu r}{n_{1}p}\log n_{2}}\sigma_{1}(\bm{M}^{\star})\le\frac{1}{2}\sigma_{r}(\bm{M}^{\star}),
\]
as long as $n_{1}p\geq C_{1}\kappa^{2}\mu r\log n_{2}$ for some sufficiently
large constant $C_{1}>0$. We are now prepared to apply Wedin's $\sin\Theta$
theorem to obtain the following upper bounds on the singular space
variations. 

\begin{theorem}[Spectral method for matrix completion]Suppose that
$n_{1}p\geq C_{1}\kappa^{2}\mu r\log n_{2}$ for some sufficiently
large constant $C_{1}>0$. Then with probability exceeding $1-O(n^{-10})$,
\begin{align*}
\max\left\{ \mathsf{dist}\left(\bm{U},\bm{U}^{\star}\right),\mathsf{dist}\left(\bm{V},\bm{V}^{\star}\right)\right\}  & \lesssim\kappa\sqrt{\frac{\mu r}{n_{1}p}\log n_{2}};\\
\max\left\{ \mathsf{dist}_{\mathrm{F}}\left(\bm{U},\bm{U}^{\star}\right),\mathsf{dist}_{\mathrm{F}}\left(\bm{V},\bm{V}^{\star}\right)\right\}  & \lesssim\kappa\sqrt{\frac{\mu r^{2}}{n_{1}p}\log n_{2}}.
\end{align*}
\end{theorem}

\begin{theorem}[Matrix completion: recovered matrix]Suppose that
$n_{2}p\geq C\mu r\log n_{2}$ for some sufficiently large constant
$C>0$. Then with probability at least $1-O(n_{2}^{-10})$, one has
\begin{align*}
\|\bm{U}\bm{\Sigma}\bm{V}^{\top}-\bm{M}^{\star}\| & \lesssim\sqrt{\frac{\mu r}{n_{1}p}\log n_{2}}\,\sigma_{1}(\bm{M}^{\star}),\quad\text{and}\quad\|\bm{U}\bm{\Sigma}\bm{V}^{\top}-\bm{M}^{\star}\|_{\mathrm{F}}\lesssim\sqrt{\frac{\mu r^{2}}{n_{1}p}\log n_{2}}\,\sigma_{1}(\bm{M}^{\star}).
\end{align*}
\end{theorem}\begin{proof}By definition, $\bm{U}\bm{\Sigma}\bm{V}^{\top}$
is the best rank-$r$ approximation to $p^{-1}\mathcal{P}_{\Omega}(\bm{M}^{\star})$,
i.e.~$\|\bm{U}\bm{\Sigma}\bm{V}^{\top}-p^{-1}\mathcal{P}_{\Omega}(\bm{M}^{\star})\|=\min_{\bm{A}:\mathsf{rank}(\bm{A})\leq r}\|\bm{A}-p^{-1}\mathcal{P}_{\Omega}(\bm{M}^{\star})\|$.
As a consequence, 
\[
\|\bm{U}\bm{\Sigma}\bm{V}^{\top}-\bm{M}^{\star}\|\leq\|\bm{U}\bm{\Sigma}\bm{V}^{\top}-p^{-1}\mathcal{P}_{\Omega}(\bm{M}^{\star})\|+\|p^{-1}\mathcal{P}_{\Omega}(\bm{M}^{\star})-\bm{M}^{\star}\|\leq2\|p^{-1}\mathcal{P}_{\Omega}(\bm{M}^{\star})-\bm{M}^{\star}\|,
\]
where the first relation is the triangle inequality. This combined
with Lemma XXX finishes the proof of the first part. Regarding the
second part, i.e.~the bound on the Frobenius norm, we make the observation
that $\bm{U}\bm{\Sigma}\bm{V}^{\top}-\bm{M}^{\star}$ has rank at
most $2r$, which immediately implies 
\[
\|\bm{U}\bm{\Sigma}\bm{V}^{\top}-\bm{M}^{\star}\|_{\mathrm{F}}\leq\sqrt{2r}\|\bm{U}\bm{\Sigma}\bm{V}^{\top}-\bm{M}^{\star}\|\lesssim\sqrt{\frac{\mu r^{2}}{n_{1}p}\log n_{2}}\sigma_{1}(\bm{M}^{\star}).
\]
This completes the proof of the second part. \end{proof}

\bibliographystyle{plain}
\bibliography{\string"/Users/congma/Dropbox (Princeton)/Research/SpectralMethod_tutorial/FnT/bibfile_Spectral\string"}

\end{document}
